% problem_id: S046-thm:ntype-nloop
% source_id: S046 (HoTT Book)
% source_file: hlevels.tex
% statement_locator: hlevels.tex lines 423--426
% proof_locator: hlevels.tex lines 427--445
% env: thm
% pairing: tight (verified)
% status: complete (statement + author proof verbatim + reason + locators)
%
\begin{thm}\label{thm:ntype-nloop}
  \index{loop space!iterated}%
  For every $n\ge -1$, a type $A$ is an $n$-type if and only if $\Omega^{n+1}(A,a)$ is contractible for all $a:A$.
\end{thm}

\begin{proof}
  Recalling that $\Omega^0(A,a) = (A,a)$, the case $n=-1$ is \cref{ex:prop-inhabcontr}.
  The case $n=0$ is \cref{thm:h-set-uip-K}.
  Now we use induction; suppose the statement holds for $n:\N$.
  By \cref{thm:hlevel-loops}, $A$ is an $(n+1)$-type iff $\Omega(A,a)$ is an $n$-type for all $a:A$.
  By the inductive hypothesis, the latter is equivalent to saying that $\Omega^{n+1}(\Omega(A,a),p)$ is contractible for all $p:\Omega(A,a)$.

  Since $\Omega^{n+2}(A,a) \defeq \Omega^{n+1}(\Omega(A,a),\refl{a})$, and $\Omega^{n+1} = \Omega^n \circ \Omega$, it will suffice to show that $\Omega(\Omega(A,a),p)$ is equal to $\Omega(\Omega(A,a),\refl{a})$, in the type $\pointed\type$ of pointed types.
  For this, it suffices to give an equivalence
  \[ g : \Omega(\Omega(A,a),p) \eqvsym \Omega(\Omega(A,a),\refl{a}) \]
  which carries the basepoint $\refl{p}$ to the basepoint $\refl{\refl{a}}$.
  For $q:p=p$, define $g(q):\refl{a} = \refl{a}$ to be the following composite:
  \[ \refl{a} = p\ct \opp p \overset{q}{=} p\ct\opp p = \refl{a}, \]
  where the path labeled ``$q$'' is actually $\apfunc{\lam{r} r\ct\opp p} (q)$.
  Then $g$ is an equivalence because it is a composite of equivalences
  \[ (p=p) \xrightarrow{\apfunc{\lam{r} r\ct\opp p}} (p\ct \opp p = p\ct \opp p) \xrightarrow{i\ct - \ct \opp i} (\refl{a} = \refl{a}). \]
  using \cref{eg:concatequiv,thm:paths-respects-equiv}, where $i:\refl{a} = p\ct \opp p$ is the canonical equality.
  And it is evident that $g(\refl{p}) = \refl{\refl{a}}$.
\end{proof}
